\section{Galois representations of the full Hecke points}

% BEGIN H86P1
\paragraph{H86P1. An actual base-change list at every localized label.}
The host satisfies $E=FK$ for $K=\mathbb Q(\sqrt{-d})$: write
$x=(x+cx)/2+\sqrt{-d}\,(x-cx)/(2\sqrt{-d})$ with both coefficients
fixed by $c$. Its prescribed real degree is at least two, and its CM
extension is unramified at finite places by the host data. Every
place above $2$ splits. Indeed $d+1=8m$ and the conjugate integers
$(1\pm\sqrt{-d})/2$ have sum $1$ and product $2m$.
Exactly one is in a given prime above $2$, which is therefore moved
by conjugation. The member's positive definite Hermitian spaces give
definiteness for all rows $1\le j\le m_{\rm row}$ in the design range.

Now fix $P\in\mathrm{AV}_i$, including any representative with
no Hecke point. The member's base-change datum embeds it in the
unitary labels of row $j(i)$ and gives sphericity there at every
nonsplit place. The positive rank is $n=r+t$. Apply B85U2 to this
image with the field hypotheses just checked. The embedding preserves
the split local parameter, the nonsplit spherical tuple and the
infinity multiset. Consequently the \emph{same} list $\mathcal P$
satisfies all six clauses of $\mathrm{IsBCV}(i,P,\mathcal P)$.
This proves $\mathrm{adm1}$ on the selected member.
% END H86P1

% BEGIN H86P2
\paragraph{H86P2. The representation attached to a list and its weights at every coefficient place.}
Fix any list from H86P1, including an alternative list satisfying
the same tests. For each $(a,\pi,d)$ in it and $0\le k<d$, put
$e=d-2k-n$. The $k$th term of the Speh parameter after overall
normalization is $\pi|\det|^{e/2}$. If $z$ is an exponent of
$\pi$ at a complex embedding, then $z+(d-1)/2-k$ occurs in the
given infinity multiset of $P$. This multiset is regular and lies
in $(n-1)/2+\mathbb Z$ by its dominant coefficient weight.
Thus $z+e/2\in\mathbb Z$ and the exponents within the term
are distinct.

Here is the half-integral twist in the Galois construction. Write
$u=(e+a-1)/2$. If $u\in\mathbb Z$, $\pi$ is regular algebraic
in rank $a$. Take the usual attachment, normalized by $(1-a)/2$,
and tensor with the algebraic character for $|\cdot|^u$ in the
convention of S86A4. If $u\in\mathbb Z+1/2$, first replace
$\pi$ by $\pi\eta$. At each embedding the exponent of $\eta$
is $1/2$ or $-1/2$. Thus the new exponents lie in
$(a-1)/2+\mathbb Z$, still without repetition. This twist preserves
cuspidality and conjugate self-duality. Tensor its attachment with
the avatar of $\eta^{-1}|\cdot|^u$, an algebraic character since
its infinity exponents are integral. In both cases the auxiliary
character has unramified smooth part at every finite place. The
normalized finite parameter has exactly the requested twist, and
the jumps are $-(z+e/2)$. For example in the second case the sum of
the jumps of the factors is
$(a-1)/2-(z+h)-(-h+u)=-z-e/2$, with $h$ the exponent of $\eta$.

Taking the direct sum over all terms gives a continuous semisimple
representation $R_{\mathcal P}$ of dimension $n$. At each place
where the list is spherical, every term has unramified parameter
and zero monodromy: inertia and the nilpotent operator are
block-diagonal on the sum. Away from $\ell$ its parameter is the
normalized list at these places, and above $\ell$ it is crystalline
there by the spherical case of A72F8 or A72F3. At every $w\mid\ell$,
including the two type primes when applicable, every factor is de
Rham by the full attachment theorem and the rank-one comparison.
For $\tau$ inducing $w$, the jumps are
\[
        \bigl\{(n-1)/2-z':z'\in\mathrm{infWt}(i,\iota\tau)\bigr\}
\]
with multiplicity. In particular these jumps are distinct at every
embedding. This application uses the full de Rham assertion for
regular algebraic cusps even at places where the source asks only
for a conditional Weil--Deligne comparison.
% END H86P2

% BEGIN H86P3
\paragraph{H86P3. Existence and uniqueness at a point.}
For $x\in X_i$ take any $P$ above $x$ by the member's point
surjection and any base-change list by H86P1. At a split place
outside $S_i,Q_i$ and $\ell$, the original level gives a spherical
parameter. Reading the Frobenius of H86P2 gives the original Hecke
polynomial of $x$, by the split-point identity including its norm
normalization. Any lift of that Frobenius has the same reading by
unramifiedness. Thus $R_{\mathcal P}$ is a rational point
representation. Two semisimple point representations have equal
characteristic polynomials on this dense set by S86A1. Continuity
of the coefficients gives equality everywhere, and Brauer--Nesbitt
gives an isomorphism. This includes the scalar extension of every
integral good point representation: semisimplicity persists over
an extension of a characteristic-zero field, by separability of the
semisimple image algebra. Conversely the member's original large
coefficient duty (M84T6) realizes each rational point representation
up to conjugacy over $O_i$. Its matrices are continuous for the
valuation topology of this ring, the original Frobenius identities
descend by injectivity and semisimplicity descends to its fraction
field. Choosing at each point gives the required integral family.

Two labels above the same point consequently give isomorphic list
representations. At all but finitely many places away from $\ell$,
both lists are spherical, and H86P2 and cancellation of the common
unramified twist identify their parameters. S86A1 recovers the
original lists up to permutation, including the Speh lengths.
Their parameters hence agree at \emph{all} places. The two
$\mathrm{IsBCV}$ readings give the requested equal split parameters
and equal base changes of the nonsplit Satake tuples.
% END H86P3

% BEGIN H86P4
\paragraph{H86P4. The four integral family properties.}
Fix any integral good family $R_x$ at $i$. For each $j$ and each
$g\in G_E$, the $j$th coefficient of the reversed characteristic
polynomial of $R_x(g)$ is a function of $x$. This function belongs to
the original integral Hecke algebra $T_i\subset O_i^{X_i}$.
Indeed it does so on the split good Frobenius elements by the
defining Hecke polynomial, and is continuous in $g$ in the finite
product topology. The algebra is closed there: as an $O_i$-submodule
of a finite free module it is finitely generated, and is the image
of a compact finite power of $O_i$. Use S86A1's density. This is
$G1$. Residual evaluations at different points of the same $T_i$
agree. Each evaluation followed by reduction is surjective by the
constants and has maximal kernel, necessarily the unique maximal
ideal of $T_i$. Since each map fixes the residue-field constants
their induced maps coincide. This proves independence of the full
residual characteristic polynomial at every $g$.

At $w\nmid\ell$ outside $S_i\cup Q_i$ the local list of each
point is spherical, by the split level or the inert base change.
H86P2--3 give unramifiedness in the rational model. Injectivity
of $O_i$ into the coefficient field gives trivial inertia already
integrally, at every conjugate choice of place. This is $G4$.
At $w\mid\ell$, when $\ell$ is neither type prime, such a list is
also spherical: $Q_i$ excludes these places, and the chosen
auxiliary prime is distinct from $\ell$ even when $\ell$ is the
first auxiliary prime. The split level and inert reading exhaust
the possibilities. H86P2--3 give crystallinity and the stated
labelled jumps, proving $G2a$.

Finally take one finite splitting extension of $\operatorname{Frac}(O_i)$
for the representations at every $x\in X_i$. For instance, split the
semisimple envelope of their direct sum over an algebraic closure and
adjoin the finitely many coefficients of a splitting basis and its
matrix units. All simple submodules are now absolutely irreducible.
If two disjoint simple submodules of $R_x$ were isomorphic, after
scalar extension they would give a positive-dimensional repeated
direct summand. At any one embedding above $\ell$, de Rham periods
of a direct sum are the sum of the periods. A direct summand of a
de Rham representation is de Rham, and its labelled jump multiset
has positive cardinality equal to its dimension. Every one of its
jumps would occur at least twice in the ambient multiset, contrary
to H86P2. This proves the distinctness assertion also when $\ell$ is
a type prime and the source imposes its additional activity guard.
Together with residual polynomial independence this is $G3$.
H86P3 and these four assertions prove the full
$\mathrm{provisionFortyOne}$ for every original $i$.
Of the three conjuncts of the selected field $\mathrm{Conclusion1}$,
the middle one is the trivial $\mathrm{g4p}$ of S86A4. Only
$\mathrm{Conclusion42}$, the parameter assertion for nonzero
theta rows, remains for that field.
% END H86P4
